\section{Finite collar histories and good-history sampling}
\label{sec:al_finite_scan_histories}

Fix a finite domain $\Lambda\subset\mathbb Z^2$, a color $A\subset\Lambda$,
and a nonempty finite target $T\subset\mathbb Z^2$. Write
$d(x)=\min_{t\in T}\|x-t\|_\infty$ and $T_j=T+[-j,j]^2_{\mathbb Z}$.
The graph metric is that of the induced nearest-neighbor graph on $\Lambda$.
The following finite construction follows
\cite[Section~9, moving partitions]{OpenAI2026AreaLaw}.

\begin{definition}[Physical partitions and padded fill slots]
    \label{def:al_scan_physical_partition}
    \lean{TNLean.PEPS.AreaLaw.Scan.PhysicalPartition,
        TNLean.PEPS.AreaLaw.Scan.middle,
        TNLean.PEPS.AreaLaw.Scan.receiving,
        TNLean.PEPS.AreaLaw.Scan.assign,
        TNLean.PEPS.AreaLaw.Scan.initialPartition,
        TNLean.PEPS.AreaLaw.Scan.orientedDepth,
        TNLean.PEPS.AreaLaw.Scan.initialFront,
        TNLean.PEPS.AreaLaw.Scan.fillCount,
        TNLean.PEPS.AreaLaw.Scan.fillSide,
        TNLean.PEPS.AreaLaw.Scan.nominalFront,
        TNLean.PEPS.AreaLaw.Scan.depthRow,
        TNLean.PEPS.AreaLaw.Scan.orientedRow,
        TNLean.PEPS.AreaLaw.Scan.fillSlot,
        TNLean.PEPS.AreaLaw.Scan.fill,
        TNLean.PEPS.AreaLaw.Scan.charge}
    \leanok
    A physical partition assigns each site to the near side $P_0$, the middle
    $Y$, or the far side $V$. For cutoffs $u<v$, initially
    $P_0=\{x\in A:d(x)\le u\}$ and
    $V=(\Lambda\setminus A)\cup\{x\in A:d(x)>v\}$.
    Order each positive-depth row once and pad it to $n$ slots. Sites outside
    $A$ and absent entries are blank. The far side visits rows in decreasing
    depth, using the same fixed order within each row.
    Fills alternate near, far, starting near. After $k$ fills the two consumed
    slot counts are $t_P=\lfloor(k+1)/2\rfloor$ and $t_F=\lfloor k/2\rfloor$;
    the oriented fronts are
    $j_P=u+1+\lfloor t_P/n\rfloor$ and $j_F=-v+\lfloor t_F/n\rfloor$.
    A move of a set $B$ transfers only $B\cap Y$ to its receiving side.
    A charge transfers this set only if $B$ meets both that side and $Y$.
\end{definition}

\begin{theorem}[Consumed slots and the actual middle]
    \label{thm:al_scan_consumed_slots}
    \lean{TNLean.PEPS.AreaLaw.Scan.assign_eq_none,
        TNLean.PEPS.AreaLaw.Scan.middle_assign_subset,
        TNLean.PEPS.AreaLaw.Scan.middle_sdiff_middle_assign,
        TNLean.PEPS.AreaLaw.Scan.assign_of_assigned,
        TNLean.PEPS.AreaLaw.Scan.fillCount_succ,
        TNLean.PEPS.AreaLaw.Scan.mem_depthRow,
        TNLean.PEPS.AreaLaw.Scan.mem_orientedRow,
        TNLean.PEPS.AreaLaw.Scan.fill_blank,
        TNLean.PEPS.AreaLaw.Scan.fill_assigned,
        TNLean.PEPS.AreaLaw.Scan.charge_moves_middle_part,
        TNLean.PEPS.AreaLaw.Scan.exists_fillSlot,
        TNLean.PEPS.AreaLaw.Scan.fill_at_side_index,
        TNLean.PEPS.AreaLaw.Scan.side_index_lt,
        TNLean.PEPS.AreaLaw.Scan.initial_middle_depth,
        TNLean.PEPS.AreaLaw.Scan.nominalFront_le_of_unconsumed}
    \leanok
    \uses{def:al_scan_physical_partition}
    A blank or already assigned fill slot makes no physical move but increases
    its receiving side's consumed count by one. Assignments never change sides,
    and the middle only decreases. For a selected set $B$, the removed middle
    is exactly $B\cap Y$. If all rows have at most $n$ sites, a site remaining
    in the middle after $k$ fills has oriented depth at least the corresponding
    nominal front.
\end{theorem}
\begin{proof}\leanok
    If a site's row is completed, its fixed index is less than $n$, so it occurs
    in a consumed slot. A site initially outside the middle never enters it,
    and consuming its slot removes it from the middle regardless of earlier
    charges. Near slot $t$ is consumed at fill $2t$, far slot $t$ at fill $2t+1$.
\end{proof}

\begin{definition}[Finite histories and their physical partitions]
    \label{def:al_scan_actual_histories}
    \lean{TNLean.PEPS.AreaLaw.Scan.CollarScan,
        TNLean.PEPS.AreaLaw.Scan.CollarScan.M,
        TNLean.PEPS.AreaLaw.Scan.CollarScan.lower,
        TNLean.PEPS.AreaLaw.Scan.CollarScan.upper,
        TNLean.PEPS.AreaLaw.Scan.CollarScan.ball,
        TNLean.PEPS.AreaLaw.Scan.CollarScan.front,
        TNLean.PEPS.AreaLaw.Scan.CollarScan.candidates,
        TNLean.PEPS.AreaLaw.Scan.CollarScan.selected,
        TNLean.PEPS.AreaLaw.Scan.CollarScan.chargeStep,
        TNLean.PEPS.AreaLaw.Scan.CollarScan.pairStep,
        TNLean.PEPS.AreaLaw.Scan.CollarScan.bandState,
        TNLean.PEPS.AreaLaw.Scan.CollarScan.state,
        TNLean.PEPS.AreaLaw.Scan.CollarScan.oldChargeState,
        TNLean.PEPS.AreaLaw.Scan.ChargeChoices,
        TNLean.PEPS.AreaLaw.Scan.History,
        TNLean.PEPS.AreaLaw.Scan.extendHistory,
        TNLean.PEPS.AreaLaw.Scan.chargeWeight,
        TNLean.PEPS.AreaLaw.Scan.historyWeight}
    \leanok
    \uses{def:al_scan_physical_partition}
    For each band $0\le g<K$, choose an independent uniform offset
    $r_g\in\{0,\ldots,m-1\}$ and set
    $u_g=8gm+m+r_g$, $v_g=8gm+5m+r_g$.
    Each charge chooses, independently in every band, a uniform side and a
    uniform slot among $M$ fixed slots. A history after $k$ fill--charge pairs
    consists of these $K$ offsets and $kK$ side--slot choices.
    Its classical weight is $m^{-K}(2M)^{-Kk}$; a new charge has weight
    $(2M)^{-K}$. Evaluate each history recursively: first perform the next
    prescribed fill, then the chosen charge at the resulting front.
    A nonblank labelled slot selects its graph $r_0$-ball; a blank slot makes
    no move.

    The picture shows the nominal partitions of $A$ in two consecutive bands,
    ordered by depth, with $X=A\cap T$ and $U=P_0\setminus X$; arrows show the
    fill directions at the fronts, and brackets mark the charge intervals of
    oriented depth $[j-r_0,j+D]$ on the two sides.
    \begin{tenkzequation}
        % Source: arXiv OpenAI2026AreaLaw, sections/08-scanner.tex:156--192
        % (figure scanner:figure), prose lines 92--130; fronts (scanner:fronts).
        % Geometry only: there are no tensors, contractions or open indices,
        % and the boundary signature is empty.
        % Horizontal axis: unsigned depth d(x), increasing to the east
        % (column 2 is d=0); it is not a lattice direction. Not to scale:
        % r_0 << D << m and the band spacing 8m are compressed.
        % Rows 2 and 5: bands g and g+1; each dot is one depth row of A
        % (a dot is a set of sites, not a tensor). Tinted enclosures are the nominal
        % parts X=A cap T, U=P_0 minus X, the middle Y and the far side V
        % (only sites of A are drawn; Lambda minus A also lies in V).
        % Rows 1 and 4: fill arrows. The near front j_P (first column of Y)
        % fills toward larger depth; the far front, at unsigned depth -j_F
        % (last column of Y, j_F being oriented), fills toward smaller depth.
        % Brackets under rows 2 and 5: charge intervals. Near side: depths
        % [j_P-r_0, j_P+D]; far side: oriented depths [j_F-r_0, j_F+D], i.e.
        % unsigned depths [-j_F-D, -j_F+r_0]; drawn with r_0=D=1 column.
        % Charges can assign sites inside these intervals, so the actual U,
        % Y, V need not be depth intervals.
        \tenkzkernel
        \begin{tenkz}[size=s, metrics=compact, rows={wire,wire,wire,wire,wire,wire,wire},
                cols=18, bonds=none]
            \tndeclare{species}{target}{hue=source:blue}
            \tndeclare{species}{near}{hue=source:green!50!black}
            \tndeclare{species}{middle}{hue=source:orange!85!black}
            \tndeclare{species}{far}{hue=source:gray}
            \tndeclare{species}{charge}{hue=source:black}
            \tnmark[form=label, label pos=w]{(2,1)}{$g$}
            \tn[at=(2,2)]{} \tn[at=(2,3)]{} \tn[at=(2,4)]{} \tn[at=(2,5)]{}
            \tn[at=(2,6)]{} \tn[at=(2,7)]{} \tn[at=(2,8)]{} \tn[at=(2,9)]{}
            \tn[at=(2,10)]{} \tn[at=(2,11)]{} \tn[at=(2,12)]{} \tn[at=(2,13)]{}
            \tn[at=(2,14)]{} \tn[at=(2,15)]{} \tn[at=(2,16)]{} \tn[at=(2,17)]{}
            \tn[at=(2,18)]{}
            \tnmark[form=enclosure, species=target, tint, label pos=n]{(2,2)}{$X$}
            \tnmark[form=enclosure, species=near, tint, label pos=n]{(2,3) .. (2,5)}{$U$}
            \tnmark[form=enclosure, species=middle, tint, label pos=n]{(2,6) .. (2,11)}{$Y$}
            \tnmark[form=enclosure, species=far, tint, label pos=n]{(2,12) .. (2,18)}{$V$}
            \tnwire[kind=string, dir=to, name=scanNearFrontG]{(1,6)}{(1,7)}
            \tnmark[form=label, label pos=n]{on scanNearFrontG 0.5}{$j_P$}
            \tnwire[kind=string, dir=to, name=scanFarFrontG]{(1,11)}{(1,10)}
            \tnmark[form=label, label pos=n]{on scanFarFrontG 0.5}{$-j_F$}
            \tnmark[form=bracket, species=charge, label pos=s]{(2,5) .. (2,7)}{}
            \tnmark[form=bracket, species=charge, label pos=s]{(2,10) .. (2,12)}{}
            \tnmark[form=label, label pos=w]{(5,1)}{$g+1$}
            \tn[at=(5,2)]{} \tn[at=(5,3)]{} \tn[at=(5,4)]{} \tn[at=(5,5)]{}
            \tn[at=(5,6)]{} \tn[at=(5,7)]{} \tn[at=(5,8)]{} \tn[at=(5,9)]{}
            \tn[at=(5,10)]{} \tn[at=(5,11)]{} \tn[at=(5,12)]{} \tn[at=(5,13)]{}
            \tn[at=(5,14)]{} \tn[at=(5,15)]{} \tn[at=(5,16)]{} \tn[at=(5,17)]{}
            \tn[at=(5,18)]{}
            \tnmark[form=enclosure, species=target, tint, label pos=n]{(5,2)}{$X$}
            \tnmark[form=enclosure, species=near, tint, label pos=n]{(5,3) .. (5,11)}{$U$}
            \tnmark[form=enclosure, species=middle, tint, label pos=n]{(5,12) .. (5,17)}{$Y$}
            \tnmark[form=enclosure, species=far, tint, label pos=n]{(5,18)}{$V$}
            \tnwire[kind=string, dir=to, name=scanNearFrontH]{(4,12)}{(4,13)}
            \tnmark[form=label, label pos=n]{on scanNearFrontH 0.5}{$j_P$}
            \tnwire[kind=string, dir=to, name=scanFarFrontH]{(4,17)}{(4,16)}
            \tnmark[form=label, label pos=n]{on scanFarFrontH 0.5}{$-j_F$}
            \tnmark[form=bracket, species=charge, label pos=s]{(5,11) .. (5,13)}{}
            \tnmark[form=bracket, species=charge, label pos=s]{(5,16) .. (5,18)}{}
            \tnwire[kind=string, dir=to, name=scanDepthAxis]{(7,2)}{(7,18)}
            \tnmark[form=label, label pos=e]{(7,18)}{$d$}
        \end{tenkz}
    \end{tenkzequation}
\end{definition}

\begin{theorem}[Normalized classical histories and deterministic fronts]
    \label{thm:al_scan_history_weights}
    \lean{TNLean.PEPS.AreaLaw.Scan.card_history,
        TNLean.PEPS.AreaLaw.Scan.chargeWeight_nonneg,
        TNLean.PEPS.AreaLaw.Scan.historyWeight_nonneg,
        TNLean.PEPS.AreaLaw.Scan.sum_chargeWeight,
        TNLean.PEPS.AreaLaw.Scan.sum_historyWeight,
        TNLean.PEPS.AreaLaw.Scan.historyWeight_extendHistory,
        TNLean.PEPS.AreaLaw.Scan.sum_historyWeight_extendHistory,
        TNLean.PEPS.AreaLaw.Scan.sum_chargeWeight_band,
        TNLean.PEPS.AreaLaw.Scan.sum_historyWeight_extendHistory_band,
        TNLean.PEPS.AreaLaw.Scan.sum_historyWeight_offset,
        TNLean.PEPS.AreaLaw.Scan.CollarScan.state_extendHistory,
        TNLean.PEPS.AreaLaw.Scan.CollarScan.front_extendHistory,
        TNLean.PEPS.AreaLaw.Scan.CollarScan.chargeStep_moves_middle_part,
        TNLean.PEPS.AreaLaw.Scan.CollarScan.bandState_unassigned,
        TNLean.PEPS.AreaLaw.Scan.CollarScan.oldChargeState_initial_none,
        TNLean.PEPS.AreaLaw.Scan.CollarScan.oldChargeState_front_le}
    \leanok
    \uses{def:al_scan_actual_histories,thm:al_scan_consumed_slots}
    For $m,M>0$, the history weights are nonnegative and sum to one.
    Extending a history multiplies its weight by $(2M)^{-K}$.
    Each band's offset has marginal probability $1/m$, even after subsequent
    charges; each side--slot pair at the next charge has conditional probability
    $1/(2M)$. Charges leave the nominal fronts unchanged.
    At every old charge history, each unassigned site lies on or beyond both
    oriented nominal fronts, provided its row has at most $n$ sites.
\end{theorem}
\begin{proof}\leanok
    There are $m^K(2M)^{Kk}$ histories and $(2M)^K$ possible new charges.
    Fixing one coordinate divides the corresponding count by $m$ or $2M$.
    Induction on the recorded choices shows that every surviving middle site
    was initially unassigned and avoids every consumed fill slot. Apply
    Theorem~\ref{thm:al_scan_consumed_slots}.
\end{proof}

\begin{theorem}[Ambient rows and graph-ball localization]
    \label{thm:al_scan_ambient_rows}
    \lean{TNLean.PEPS.AreaLaw.Scan.ambientSupDistance,
        TNLean.PEPS.AreaLaw.Scan.ambientDepth,
        TNLean.PEPS.AreaLaw.Scan.ambientDepth_nonneg,
        TNLean.PEPS.AreaLaw.Scan.abs_ambientDepth_sub_le,
        TNLean.PEPS.AreaLaw.Scan.ambientDepth_le_iff_mem_dilation,
        TNLean.PEPS.AreaLaw.Scan.ambientDepth_eq_iff_mem_dilation_sdiff,
        TNLean.PEPS.AreaLaw.Scan.abs_ambientDepth_sub_le_walk_length,
        TNLean.PEPS.AreaLaw.Scan.abs_ambientDepth_sub_le_of_edist_le,
        TNLean.PEPS.AreaLaw.Scan.ambientSupDistance_le_one_of_domainGraph_adj,
        TNLean.PEPS.AreaLaw.Scan.abs_ambientDepth_sub_le_domainGraph,
        TNLean.PEPS.AreaLaw.Scan.depthRow_ambientDepth_length_le,
        TNLean.PEPS.AreaLaw.Scan.depthRow_ambientDepth_length_le_of_row_bound,
        TNLean.PEPS.AreaLaw.Scan.domainGraph_ball_subset_of_clearance}
    \leanok
    For positive $d$, the ambient depth row is $T_d\setminus T_{d-1}$;
    its physical part has no larger cardinality. Graph distance bounds depth
    variation: $|d(x)-d(y)|\le d_\Lambda(x,y)$.
    Let $Z$ be the actual endpoints of edges crossing $A$.
    If $\operatorname{dist}_\infty(T,Z)>2L+10r_0$, every graph $r_0$-ball
    meeting $A\cap T_L$ is wholly contained in $A$.
\end{theorem}
\begin{proof}\leanok
    Minimize the integer sup-norm distance over $T$.
    Each lattice edge changes that minimum by at most one; sum along a
    shortest graph walk. If a ball meeting $A\cap T_L$ leaves $A$, a path
    of length at most $2r_0$ within the two radii crosses an edge with an
    endpoint of depth at most $L+2r_0$, contradicting the clearance.
\end{proof}

\begin{theorem}[Uniform labelled charge capacity]
    \label{thm:al_scan_charge_capacity}
    \lean{TNLean.PEPS.AreaLaw.Scan.chargeCandidates,
        TNLean.PEPS.AreaLaw.Scan.mem_chargeCandidates,
        TNLean.PEPS.AreaLaw.Scan.card_chargeCandidates_le,
        TNLean.PEPS.AreaLaw.Scan.chargeSlotCount,
        TNLean.PEPS.AreaLaw.Scan.card_chargeCandidates_le_chargeSlotCount,
        TNLean.PEPS.AreaLaw.Scan.orderedChargeCandidates,
        TNLean.PEPS.AreaLaw.Scan.paddedChargeSlot,
        TNLean.PEPS.AreaLaw.Scan.mem_of_paddedChargeSlot_eq_some,
        TNLean.PEPS.AreaLaw.Scan.paddedChargeSlot_eq_none_iff,
        TNLean.PEPS.AreaLaw.Scan.existsUnique_paddedChargeSlot,
        TNLean.PEPS.AreaLaw.Scan.existsUnique_orderedChargeSlot,
        TNLean.PEPS.AreaLaw.Scan.chargeSlotCount_pos,
        TNLean.PEPS.AreaLaw.Scan.chargeSlotCount_le,
        TNLean.PEPS.AreaLaw.Scan.chargeSlotWeight_lower_bound,
        TNLean.PEPS.AreaLaw.Scan.fillCount_div_le_window,
        TNLean.PEPS.AreaLaw.Scan.bandChargeInterval_bounds,
        TNLean.PEPS.AreaLaw.Scan.bandChargeInterval_row_bound,
        TNLean.PEPS.AreaLaw.Scan.card_bandChargeCandidates_le_chargeSlotCount}
    \leanok
    \uses{def:al_scan_actual_histories,thm:al_scan_ambient_rows}
    Suppose $n,m\ge2$, $D\ge1$, $r_0\le D$, $4D\le m$, $8Km\le L$, and at
    most $nm$ fills are consumed. Assume $|T_d\setminus T_{d-1}|\le n$ for
    $1\le d\le L$, and at most $\mu$ labels have any fixed anchor.
    Each charge interval $[j-r_0,j+D]$ visits only unsigned depths in $[1,L]$.
    Its number of labels is at most $(D+r_0+1)n\mu\le3\mu nD$.
    Thus $C_1\ge3\mu$ permits padding every list to $M=\lceil C_1nD\rceil$,
    with a unique slot for every label, including labels with equal anchors.
    If $C_1>0$, then $M\le(C_1+1)nD$ and
    $1/(2M)\ge[2(C_1+1)]^{-1}/(nD)$.
\end{theorem}
\begin{proof}\leanok
    Sum the row cardinalities, then sum the anchor fibers. Either nominal front
    advances at most $m$ rows; the band cutoffs and $4D\le m$ keep the entire
    candidate interval in the collar. Ordering distinct labels gives distinct
    slots. Finally $\lceil C_1nD\rceil\le C_1nD+1\le(C_1+1)nD$.
\end{proof}

\begin{theorem}[Actual ball sampling on a good history]
    \label{thm:al_scan_good_ball_sampling}
    \lean{TNLean.PEPS.AreaLaw.Scan.CollarScan.IsGoodOldHistory,
        TNLean.PEPS.AreaLaw.Scan.CollarScan.good_split_anchor_bounds,
        TNLean.PEPS.AreaLaw.Scan.CollarScan.existsUnique_selected_of_bounds,
        TNLean.PEPS.AreaLaw.Scan.CollarScan.selected_split_sampling,
        TNLean.PEPS.AreaLaw.Scan.CollarScan.oldChargeState_depth_bounds,
        TNLean.PEPS.AreaLaw.Scan.CollarScan.good_split_sampling_domainGraph}
    \leanok
    \uses{thm:al_scan_history_weights,thm:al_scan_ambient_rows,thm:al_scan_charge_capacity}
    Retain the capacity and clearance assumptions, and suppose $2r_0\le D$.
    An old charge history is good when every assigned site \emph{of $A$}
    has oriented depth at most $j+D/2$ on its receiving side.
    For every labelled graph $r_0$-ball meeting that side and the actual middle,
    the anchor lies in $[j-r_0,j+D]$. Its labelled slot is unique.
    Choosing this side and slot has probability exactly $1/(2M)$, at least
    $c/(nD)$ for $c=1/[2(C_1+1)]>0$, and the extended physical history removes
    exactly the ball's old $Y$-part from $Y$.
\end{theorem}
\begin{proof}\leanok
    An unassigned site has oriented depth at least $j$, so the anchor has depth
    at least $j-r_0$. The ball lies in $A$ by clearance; an assigned site of
    the ball therefore has depth at most $j+D/2$, giving anchor depth at most
    $j+D/2+r_0\le j+D$. The unique-slot and classical marginal formulas give
    the probability. The actual charge transition transfers precisely the
    selected ball intersected with the old middle.
\end{proof}

\begin{theorem}[Actual sampling of truncated designated supports]
    \label{thm:al_scan_designated_sampling}
    \lean{TNLean.PEPS.AreaLaw.Scan.CollarScan.truncationSet,
        TNLean.PEPS.AreaLaw.Scan.CollarScan.middle_subset_truncationSet,
        TNLean.PEPS.AreaLaw.Scan.CollarScan.designatedSupport_eq_ball_of_split,
        TNLean.PEPS.AreaLaw.Scan.CollarScan.good_designated_split_sampling_domainGraph}
    \leanok
    \uses{thm:al_crossing_budget,thm:al_scan_good_ball_sampling}
    Set $S_0=A\cap T_L$ and retain the good-history sampling assumptions.
    Every truncated designated support meeting a receiving side and the old
    middle is a graph $r_0$-ball. Its labelled side--slot event has probability
    exactly $1/(2M)\ge c/(nD)$, and the actual extended history removes exactly
    the designated support's old unassigned part.
\end{theorem}
\begin{proof}\leanok
    The actual old middle is contained in $S_0$. The receiving-side witness
    lies outside that middle, so the designated support crosses a subset of
    $S_0$. Theorem~\ref{thm:al_crossing_budget} identifies its truncation radius
    with $r_0$, including the potentially variable-radius and disconnected
    cases. Apply Theorem~\ref{thm:al_scan_good_ball_sampling} to this equal ball.
\end{proof}

\begin{theorem}[Expected cost of the actual charge moves]
    \label{thm:al_scan_expected_move_cost}
    \lean{TNLean.PEPS.AreaLaw.Scan.CollarScan.splitIncidences,
        TNLean.PEPS.AreaLaw.Scan.CollarScan.mem_splitIncidences,
        TNLean.PEPS.AreaLaw.Scan.CollarScan.chargeMoveCost,
        TNLean.PEPS.AreaLaw.Scan.CollarScan.splitIncidenceCost,
        TNLean.PEPS.AreaLaw.Scan.CollarScan.splitIncidenceCost_band_le_expected_chargeMoveCost,
        TNLean.PEPS.AreaLaw.Scan.CollarScan.splitIncidenceCost_le_expected_chargeMoveCost,
        TNLean.PEPS.AreaLaw.Scan.CollarScan.good_splitIncidenceCost_le_expected_chargeMoveCost_domainGraph}
    \leanok
    \uses{thm:al_scan_good_ball_sampling}
    Under the good-history sampling assumptions, let $F_{g,s}(B)\ge0$ be a
    nonnegative cost of a finite set of physical sites moved toward side $s$
    in band $g$.
    Sum $F_{g,s}(B_i\cap Y_g)$ over every splitting incidence $(g,s,i)$,
    where $s$ is the receiving side and $B_i$ is the labelled charge ball.
    If $\Delta_g(c)$ is the actual set removed from $Y_g$ by the simultaneous
    charge choice $c$ and $s_g(c)$ its selected side, then
    \begin{align}
        \frac{1}{2M}\sum_{(g,s,i)\ \mathrm{split}} F_{g,s}(B_i\cap Y_g)
         & \le \sum_c q_c\sum_g F_{g,s_g(c)}(\Delta_g(c)).
        \label{eq:al_scan_expected_move_cost}
    \end{align}
    For $C_1>0$, the left coefficient may be replaced by
    $1/[2(C_1+1)nD]$. Labels sharing an anchor retain separate summands.
\end{theorem}
\begin{proof}\leanok
    For each band and realized side--slot pair, at most one labelled incidence
    is selected. Its cost is the cost of the actual removed middle set.
    All other event contributions are zero, and the realized cost is
    nonnegative. Sum this pointwise inequality against the classical choice
    weights, use the $1/(2M)$ marginal for each available labelled slot,
    and interchange the finite sums over choices and bands.
\end{proof}

\begin{theorem}[Expected cost for the truncated designated supports]
    \label{thm:al_scan_designated_move_cost}
    \lean{TNLean.PEPS.AreaLaw.Scan.CollarScan.designatedSplitIncidences,
        TNLean.PEPS.AreaLaw.Scan.CollarScan.mem_designatedSplitIncidences,
        TNLean.PEPS.AreaLaw.Scan.CollarScan.designatedSplitIncidences_subset_splitIncidences,
        TNLean.PEPS.AreaLaw.Scan.CollarScan.designatedSplitIncidenceCost,
        TNLean.PEPS.AreaLaw.Scan.CollarScan.designatedSplitIncidenceCost_le_splitIncidenceCost,
        TNLean.PEPS.AreaLaw.Scan.CollarScan.good_designatedSplitIncidenceCost_le_expected_chargeMoveCost_domainGraph}
    \leanok
    \uses{thm:al_scan_designated_sampling,thm:al_scan_expected_move_cost}
    The expected-cost lower bound also holds with the sum restricted to
    splitting incidences of the truncated designated supports, using their
    actual old middle parts. The cost may depend on the band and receiving
    side, and repeated labels remain distinct.
\end{theorem}
\begin{proof}\leanok
    Every splitting designated support equals its charge ball, so its
    incidence and cost occur in the ball-incidence sum. All costs are
    nonnegative. Restrict that sum and apply
    Theorem~\ref{thm:al_scan_expected_move_cost}.
\end{proof}

% The actual designated-support composition is included above.
% The bad-history ancestry tail, cross-band commutation, and complete
% construction of ScanData remain separate.


\begin{theorem}[Unconditional assigned-site lead and split intervals]
    \label{thm:al_scan_unconditional_split_interval}
    \lean{TNLean.PEPS.AreaLaw.Scan.fillSlot_depth,
        TNLean.PEPS.AreaLaw.Scan.CollarScan.front_mono_succ,
        TNLean.PEPS.AreaLaw.Scan.CollarScan.selected_anchor_bounds,
        TNLean.PEPS.AreaLaw.Scan.CollarScan.bandState_assigned_lead,
        TNLean.PEPS.AreaLaw.Scan.CollarScan.oldChargeState_assigned_lead,
        TNLean.PEPS.AreaLaw.Scan.CollarScan.bandState_front_le,
        TNLean.PEPS.AreaLaw.Scan.CollarScan.old_split_anchor_bounds,
        TNLean.PEPS.AreaLaw.Scan.CollarScan.state_split_anchor_bounds,
        TNLean.PEPS.AreaLaw.Scan.CollarScan.state_depth_bounds,
        TNLean.PEPS.AreaLaw.Scan.CollarScan.old_split_anchor_bounds_domainGraph,
        TNLean.PEPS.AreaLaw.Scan.CollarScan.state_split_anchor_bounds_domainGraph}
    \leanok
    \uses{def:al_scan_actual_histories,thm:al_scan_consumed_slots,thm:al_scan_ambient_rows}
    For every finite actual history and every assigned site of $A$, its oriented
    depth is at most $j+D+r_0$, where $j$ is its side's current nominal front.
    This holds both for completed histories and immediately before a charge.
    Assume the ambient row bound through $L$, $8Km\le L$, and the stated
    clearance of $T$ from the physical cut endpoints. If a radius-$r_0$ ball
    meets a receiving side and the middle, its anchor's oriented depth $d_a$ satisfies
    \begin{align}
        j-r_0\le d_a\le j+D+2r_0.
    \end{align}
    No good-history assumption is imposed. The exterior $\Lambda\setminus A$
    is not subject to the assigned-site lead bound.
\end{theorem}
\begin{proof}\leanok
    Induct on the actual evaluator. A fill assigns only its scheduled row;
    a selected charge anchor is at most $j+D$, and its ball adds at most $r_0$.
    Later fronts do not decrease. An old middle site lies beyond the completed
    rows, giving the lower split-anchor bound. An assigned site in the ball gives
    the upper bound. The middle remains in the original compact collar, so cut
    clearance places the entire ball in $A$; ambient depth variation follows
    from induced graph distance. These conclusions hold for all finite history
    lengths, hence throughout the manuscript's finite scan horizon.
\end{proof}

\begin{theorem}[Finite uniform-offset counting]
    \label{thm:al_scan_offset_count}
    \lean{TNLean.PEPS.AreaLaw.Scan.sum_historyWeight_offset_mem,
        TNLean.PEPS.AreaLaw.Scan.sum_historyWeight_mul_eq_sum_extendHistory,
        TNLean.PEPS.AreaLaw.Scan.CollarScan.front_injective_offset,
        TNLean.PEPS.AreaLaw.Scan.CollarScan.card_front_interval_le,
        TNLean.PEPS.AreaLaw.Scan.CollarScan.sum_historyWeight_event_le,
        TNLean.PEPS.AreaLaw.Scan.CollarScan.sum_historyWeight_exists_band_side_le,
        TNLean.PEPS.AreaLaw.Scan.CollarScan.historyWeight_mixture_event_le,
        TNLean.PEPS.AreaLaw.Scan.CollarScan.split_interval_band_unique}
    \leanok
    \uses{def:al_scan_actual_histories,thm:al_scan_history_weights}
    At fixed band, side and front time, at most $a+b+1$ offsets place a fixed
    depth in $[j-a,j+b]$. Every history event contained in that interval event
    therefore has probability at most $(a+b+1)/m$, regardless of its dependence
    on later choices. If $a+b\le m$, the centers' spacing $8m$ allows at most
    one band on each side to reach the fixed depth, even across different offsets.
    The union over both sides and all bands has probability at most
    $2(a+b+1)/m$. Common convex mixtures preserve the bound.
\end{theorem}
\begin{proof}\leanok
    The front is affine with coefficient $1$ or $-1$ in its offset. Inject it
    into the integer interval $[d-b,d+a]$. The exact marginal of an offset set
    is its cardinality divided by $m$. Reindexing a new history into its old
    prefix and last charge gives exactly the leaf weights $w_hq_c$; this
    identity holds for every real-valued function of the new history. For a fixed side the possible band is
    deterministic, because two different band centers cannot both be reached
    after allowing all offsets. Sum the two side bounds. This uses marginal
    domination, not independence of a state-dependent splitting event.
\end{proof}

\begin{definition}[Designated receiving--middle incidence]
    \label{def:al_scan_designated_incidence}
    \lean{TNLean.PEPS.AreaLaw.Scan.CollarScan.designatedSplitIncidence}
    \leanok
    \uses{def:al_scan_actual_histories,thm:al_scan_designated_sampling}
    A label has a receiving--middle incidence if its designated truncated support
    meets a receiving side and the actual middle of one band. The terminal
    split event of \cite[Section~7, terminal split rule]{OpenAI2026AreaLaw}
    additionally requires exactly one receiving side and containment in a single
    part of every other band, so it is contained in this incidence event.
    These are the physical statuses $P_0,Y,V$ underlying $P=CP_0,Y,F=RV$.
    No dilution is asserted for arbitrary entropy cuts such as $U=P_0\setminus X$,
    whose fixed inner boundary is not a terminal metric split.
\end{definition}

\begin{theorem}[Actual designated-support dilution]
    \label{thm:al_scan_designated_dilution}
    \lean{TNLean.PEPS.AreaLaw.Scan.CollarScan.designatedSupport_eq_ball_of_state_split,
        TNLean.PEPS.AreaLaw.Scan.CollarScan.old_designated_split_dilution_domainGraph,
        TNLean.PEPS.AreaLaw.Scan.CollarScan.state_designated_split_dilution_domainGraph,
        TNLean.PEPS.AreaLaw.Scan.CollarScan.designated_split_mixture_dilution_domainGraph,
        TNLean.PEPS.AreaLaw.Scan.CollarScan.fill_split_mixture_dilution_domainGraph}
    \leanok
    \uses{thm:al_scan_unconditional_split_interval,thm:al_scan_offset_count,
        def:al_scan_designated_incidence,thm:al_scan_designated_sampling}
    Suppose $n,m,D>0$, $C_1>0$, $r_0\le D$, $4D\le m$, $8Km\le L$, and the
    actual ambient row and cut-clearance hypotheses hold. For every label, the
    total classical weight of designated receiving--middle incidences in any
    band is at most
    \begin{align}
        2(D+3r_0+1)/m\le 10D/m.
    \end{align}
    This bound holds for the old and new distributions of every fill or charge
    round, and for every common mixture parameter $p\in[0,1]$.
\end{theorem}
\begin{proof}\leanok
    The actual middle lies in the compact truncation set. The accepted
    variable-radius truncation result thus identifies every incident designated
    support with its radius-$r_0$ ball. Apply the preceding interval theorem with
    $a=r_0$ and $b=D+2r_0$; their sum is at most $m$. Apply the finite offset
    union bound and then convexity. The estimate concerns classical history
    weights before quantum state transport. It does not construct the remaining
    scanner data fields or the bad-history probability bound.
\end{proof}
